\documentclass[12pt]{article}

\usepackage[margin=1in]{geometry}
\usepackage{graphicx}
\usepackage{booktabs}
\usepackage{amsmath}

\title{The Effect of Paid Search Advertising on Revenue: \\
       A Replication of the eBay Natural Experiment}
\author{Jai Correa}
\date{\today}

\begin{document}

\maketitle

\section{Introduction}

The research question was trying to find the causal effectiveness of 
paid search ads on eBay's revenue.  
This matters because while firms spend billions on SEM, measuring its causal
effect is difficult. A natural experiment is used to answer the question, where
eBay turned off ads in some markets. 

\section{Experimental Setting}

Blake et al. (2014) convinced eBay to stop bidding on AdWords. They did this
by doing the following. Firstly, this was not a fully randomized trial --
the largest DMAs were excluded from treatment, so treatment and control
groups differ in baseline revenue. Therefore, the paper uses differences-in-difference
(DID) rather than simple comparison of means. 65 of 210 DMAs were selected
for treatment starting May 22, 2012. Treatment lasted 8 weeks, with the 
treatment group having their search turned off, and the control group having
their search kept on.  

\section{Data and Figures}

The following is daily revenue for 210 DMAs from April to July 2012.
The variables include date, DMA identifier, treatment period, group,
and revenue.

\begin{figure}[h]
	\centering
	\includegraphics[width=0.85\textwidth]{../output/figures/figure_5_2.png}
	\caption{Average revenue for treatment and control DMAs over time.
		 The dashed line marks May 22, 2012, when SEM was turned off
		 for the treatment group.}
	\label{fig:revenue}
\end{figure}

While both groups show similar weekly oscillation patterns, it is the 
control group has higher average revenue (the largest DMAs have been
excluded from treatment). Additionally, it is also hard to see a clear 
divergence visually at this scale.

\begin{figure}[h]
	\centering
	\includegraphics[width=0.85\textwidth]{../output/figures/figure_5_3.png}
	\caption{Log-scale average revenue difference between control and treatment
		 groups over time. The gap appears to increase after May 22.}
	\label{fig:logdiff}
\end{figure}

The above figure shows the log ratio of control to treatment revenue.
The line fluctuates around a stable level before May 22 and appearing to
increase slightly after. This visual pattern motivates the formal DID test.

\section{Results}

\input{../output/tables/did_table.tex}

The point estimate $\hat{\gamma} = -0.0066$ in log scale.
In levels, this means it equals $\exp(-0.0066) = 0.9934$, showing treatment DMAs
lost about 0.66 percent of revenue when paid search was turned off. 
The 95 percent confidence interval is $[-0.0175, 0.0043]$, including zero,
implying the effect is not statistically significant at the 5 percent level.
The economic interpretation would be that turning paid search off had a very
small (and statistically insignificant) negative effect on revenue. This
suggests that for a well-known brand like eBay, most users would find the
site through organic search anyway.

\section{Conclusion}

The DID analysis finds a small, statistically insignificant effect. 
This replicates the main finding of Blake et al. (2014).
Limitations of this replication include the following:
the data is scaled/shifted (not real revenue figures);
the treatment was not fully randomized; we use a simple DID rather than 
regression with controls.
Overall, the result has important implications for marketing ROI measurement.

\begin{thebibliography}{9}
\bibitem{blake2014}
Blake, T., C. Nosko, and S. Tadelis (2014).
``Consumer Heterogeneity and Paid Search Effectiveness: A Large-Scale Field Experiment.''
\textit{Econometrica}, 83(1): 155--174.

\bibitem{taddy2019}
Taddy, M. (2019).

\textit{Business Data Science}. McGraw-Hill, Chapter 5.
\end{thebibliography}

\end{document}
